% Options for packages loaded elsewhere
\PassOptionsToPackage{unicode}{hyperref}
\PassOptionsToPackage{hyphens}{url}
\documentclass[
]{article}
\usepackage{xcolor}
\usepackage{amsmath,amssymb}
\setcounter{secnumdepth}{-\maxdimen} % remove section numbering
\usepackage{iftex}
\ifPDFTeX
  \usepackage[T1]{fontenc}
  \usepackage[utf8]{inputenc}
  \usepackage{textcomp} % provide euro and other symbols
\else % if luatex or xetex
  \usepackage{unicode-math} % this also loads fontspec
  \defaultfontfeatures{Scale=MatchLowercase}
  \defaultfontfeatures[\rmfamily]{Ligatures=TeX,Scale=1}
\fi
\usepackage{lmodern}
\ifPDFTeX\else
  % xetex/luatex font selection
\fi
% Use upquote if available, for straight quotes in verbatim environments
\IfFileExists{upquote.sty}{\usepackage{upquote}}{}
\IfFileExists{microtype.sty}{% use microtype if available
  \usepackage[]{microtype}
  \UseMicrotypeSet[protrusion]{basicmath} % disable protrusion for tt fonts
}{}
\makeatletter
\@ifundefined{KOMAClassName}{% if non-KOMA class
  \IfFileExists{parskip.sty}{%
    \usepackage{parskip}
  }{% else
    \setlength{\parindent}{0pt}
    \setlength{\parskip}{6pt plus 2pt minus 1pt}}
}{% if KOMA class
  \KOMAoptions{parskip=half}}
\makeatother
\usepackage{longtable,booktabs,array}
\usepackage{caption}
\captionsetup[table]{skip=6pt}
\newcounter{none} % for unnumbered tables
\usepackage{calc} % for calculating minipage widths
% Correct order of tables after \paragraph or \subparagraph
\usepackage{etoolbox}
\makeatletter
\patchcmd\longtable{\par}{\if@noskipsec\mbox{}\fi\par}{}{}
\makeatother
% Allow footnotes in longtable head/foot
\IfFileExists{footnotehyper.sty}{\usepackage{footnotehyper}}{\usepackage{footnote}}
\makesavenoteenv{longtable}
\setlength{\emergencystretch}{3em} % prevent overfull lines
\providecommand{\tightlist}{%
  \setlength{\itemsep}{0pt}\setlength{\parskip}{0pt}}
\usepackage{bookmark}
\IfFileExists{xurl.sty}{\usepackage{xurl}}{} % add URL line breaks if available
\urlstyle{same}
% fallback for those not using the hyperref driver hyperxmp:
\makeatletter
\@ifundefined{xmpquote}{\newcommand{\xmpquote}[1]{#1}}{}
\makeatother
\hypersetup{
  hidelinks,
  pdfcreator={LaTeX via pandoc}}

\author{}
\date{}

\begin{document}

\section{Jamovi Step-by-Step Lab Guide --- Session
1}\label{jamovi-step-by-step-lab-guide-session-1}

\begin{quote}
\textbf{Course}: Quantitative Data Analysis (M1 -- S7)
\textbf{Software}: \href{https://www.jamovi.org/}{Jamovi} (Version 2.3+)
\textbf{Session}: Session 1 --- Data Setup \& Univariate Descriptives
\textbf{Target File}: \texttt{data/processed/sp500\_esg\_jamovi.csv}
\textbf{Final Output}: \texttt{sp500\_esg\_session1.omv}
\end{quote}

\begin{center}\rule{0.5\linewidth}{0.5pt}\end{center}

\subsection{Step 1: Open the Dataset in
Jamovi}\label{step-1-open-the-dataset-in-jamovi}

\begin{enumerate}
\def\labelenumi{\arabic{enumi}.}
\tightlist
\item
  Launch \textbf{Jamovi} on your computer.
\item
  Click the top-left \textbf{\texttt{≡} menu} (hamburger icon).
\item
  Select \textbf{\texttt{Open}} \(\rightarrow\)
  \textbf{\texttt{Browse}}.
\item
  Navigate to your project folder:
  \texttt{Group\ Work\ S\&P\ 500/data/processed/}
\item
  Select \textbf{\texttt{sp500\_esg\_jamovi.csv}} (or
  \texttt{sp500\_esg\_jamovi.xlsx}).
\item
  Click \textbf{\texttt{Open}}.
\end{enumerate}

\begin{center}\rule{0.5\linewidth}{0.5pt}\end{center}

\subsection{Step 2: Configure Variable Measure Types (Crucial
Step!)}\label{step-2-configure-variable-measure-types-crucial-step}

\begin{quote}
{[}!CAUTION{]} As Dr.~NGUYEN Anh-Tuan warns in \textbf{Slide 38}:
\emph{``If the measure type is wrong, Jamovi will silently refuse to
offer you the right test in week 2. This is the number one cause of a
lost lab session.''}
\end{quote}

Double-click on each column header in the \textbf{Data} tab and set the
\textbf{Measure type} and \textbf{Data type}:

{\def\LTcaptype{none} % do not increment counter
\begin{longtable}[]{@{}
  >{\raggedright\arraybackslash}p{(\linewidth - 6\tabcolsep) * \real{0.1880}}
  >{\raggedright\arraybackslash}p{(\linewidth - 6\tabcolsep) * \real{0.1654}}
  >{\raggedright\arraybackslash}p{(\linewidth - 6\tabcolsep) * \real{0.0677}}
  >{\raggedright\arraybackslash}p{(\linewidth - 6\tabcolsep) * \real{0.5789}}@{}}
\toprule\noalign{}
\begin{minipage}[b]{\linewidth}\raggedright
Column Name
\end{minipage} & \begin{minipage}[b]{\linewidth}\raggedright
Measure Type in Jamovi
\end{minipage} & \begin{minipage}[b]{\linewidth}\raggedright
Data Type
\end{minipage} & \begin{minipage}[b]{\linewidth}\raggedright
Missing Values / Levels
\end{minipage} \\
\midrule\noalign{}
\endhead
\bottomrule\noalign{}
\endlastfoot
\texttt{Ticker} & \textbf{ID} & Text & Identifier (never analysed) \\
\texttt{Company} & \textbf{Nominal} & Text & Company names \\
\texttt{Sector} & \textbf{Nominal} & Text & 5 levels: \emph{Consumer,
Finance, Healthcare, Industrials \& Energy, Tech \& Comms} \\
\texttt{Total\_Revenue\_B} & \textbf{Continuous} & Decimal & None \\
\texttt{Market\_Cap\_B} & \textbf{Continuous} & Decimal & None \\
\texttt{Market\_Cap\_Quartile} & \textbf{Ordinal} & Text & 4 ordered
levels: \emph{Q1 \textless{} Q2 \textless{} Q3 \textless{} Q4} \\
\texttt{Profit\_Margin} & \textbf{Continuous} & Decimal & None \\
\texttt{ROE} & \textbf{Continuous} & Decimal & 32 missing values (leave
empty) \\
\texttt{Beta} & \textbf{Continuous} & Decimal & 4 missing values \\
\texttt{Headcount} & \textbf{Continuous} & Integer & 3 missing values \\
\texttt{Overall\_Governance\_Risk} & \textbf{Continuous} & Integer &
Scores 1 to 10 \\
\texttt{Governance\_Risk\_Level} & \textbf{Ordinal} & Text & 3 ordered
levels: \emph{Low \textless{} Medium \textless{} High} \\
\end{longtable}
}

\begin{center}\rule{0.5\linewidth}{0.5pt}\end{center}

\subsection{Step 3: Run Univariate Analysis on Categorical Variables
(Slide
40)}\label{step-3-run-univariate-analysis-on-categorical-variables-slide-40}

\begin{enumerate}
\def\labelenumi{\arabic{enumi}.}
\tightlist
\item
  Click on the \textbf{\texttt{Analyses}} tab in the top ribbon.
\item
  Click \textbf{\texttt{Exploration}} \(\rightarrow\)
  \textbf{\texttt{Descriptives}}.
\item
  In the variable selection box:

  \begin{itemize}
  \tightlist
  \item
    Move \textbf{\texttt{Sector}} and
    \textbf{\texttt{Governance\_Risk\_Level}} into the
    \textbf{\texttt{Variables}} box.
  \end{itemize}
\item
  Expand the \textbf{\texttt{Statistics}} panel below:

  \begin{itemize}
  \tightlist
  \item
    Check \textbf{\texttt{Frequency\ tables}} (for qualitative
    variables).
  \item
    Uncheck Mean/Median/SD for qualitative variables (avoids the
    ``average sector = 2.4'' trap!).
  \end{itemize}
\item
  Expand the \textbf{\texttt{Plots}} panel:

  \begin{itemize}
  \tightlist
  \item
    Check \textbf{\texttt{Bar\ plot}}.
  \end{itemize}
\item
  \textbf{Expected Output in Jamovi}:

  \begin{itemize}
  \tightlist
  \item
    \texttt{Sector}: Industrials \& Energy (148, 29.4\%), Consumer (116,
    23.1\%), Tech \& Comms (108, 21.5\%), Finance (71, 14.1\%),
    Healthcare (60, 11.9\%).
  \item
    \texttt{Governance\_Risk\_Level}: High (197, 39.7\%), Low (150,
    30.2\%), Medium (149, 30.0\%).
  \end{itemize}
\end{enumerate}

\begin{center}\rule{0.5\linewidth}{0.5pt}\end{center}

\subsection{Step 4: Run Univariate Analysis \& Normality Tests on
Continuous
Variables}\label{step-4-run-univariate-analysis-normality-tests-on-continuous-variables}

\begin{enumerate}
\def\labelenumi{\arabic{enumi}.}
\tightlist
\item
  Still in \textbf{\texttt{Analyses}} \(\rightarrow\)
  \textbf{\texttt{Exploration}} \(\rightarrow\)
  \textbf{\texttt{Descriptives}} (or create a new Descriptives block):
\item
  Move the continuous variables into \textbf{\texttt{Variables}}:

  \begin{itemize}
  \tightlist
  \item
    \texttt{Profit\_Margin}
  \item
    \texttt{ROE}
  \item
    \texttt{Market\_Cap\_B}
  \item
    \texttt{Total\_Revenue\_B}
  \item
    \texttt{Overall\_Governance\_Risk}
  \end{itemize}
\item
  Expand \textbf{\texttt{Statistics}}:

  \begin{itemize}
  \tightlist
  \item
    \textbf{Central Tendency}: Check \texttt{Mean}, \texttt{Median},
    \texttt{Mode}.
  \item
    \textbf{Dispersion}: Check \texttt{Std.\ deviation},
    \texttt{Variance}, \texttt{Minimum}, \texttt{Maximum}, \texttt{IQR}
    (Interquartile range).
  \item
    \textbf{Distribution Shape}: Check \texttt{Skewness} and
    \texttt{Kurtosis}.
  \item
    \textbf{Normality}: Check \textbf{\texttt{Shapiro-Wilk}}.
  \end{itemize}
\item
  Expand \textbf{\texttt{Plots}}:

  \begin{itemize}
  \tightlist
  \item
    Check \textbf{\texttt{Histogram}}.
  \item
    Check \textbf{\texttt{Density}}.
  \item
    Check \textbf{\texttt{Box\ plot}} (and check \texttt{Outliers}).
  \item
    Check \textbf{\texttt{Q-Q\ plot}}.
  \end{itemize}
\end{enumerate}

\begin{center}\rule{0.5\linewidth}{0.5pt}\end{center}

\subsection{\texorpdfstring{Step 5: Save and Export Your \texttt{.omv}
File (Slide 37 \&
41)}{Step 5: Save and Export Your .omv File (Slide 37 \& 41)}}\label{step-5-save-and-export-your-.omv-file-slide-37-41}

\begin{enumerate}
\def\labelenumi{\arabic{enumi}.}
\tightlist
\item
  Click the top-left \textbf{\texttt{≡} menu}.
\item
  Select \textbf{\texttt{Save\ As}} \(\rightarrow\)
  \textbf{\texttt{Browse}}.
\item
  Name your file: \textbf{\texttt{sp500\_esg\_session1.omv}}.
\item
  Save it in your project folder. \emph{(The \texttt{.omv} format
  bundles both your dataset and all interactive analysis output tables
  together into a single file ready for submission to Dr.~NGUYEN
  Anh-Tuan).}
\end{enumerate}

\end{document}
